% Source-derived chapter 04: Splittings of perverse filtrations
% Original source: hitchin-fibrations-abelian-surfaces-pw
\chapter{Splittings of perverse filtrations}
\label{hitchin-fibrations-abelian-surfaces-pw-chapter-04}
\source{hitchin-fibrations-abelian-surfaces-pw}

\begin{remark}[Source section]
\label{hitchin-fibrations-abelian-surfaces-pw-chapter-04-source}
\source{hitchin-fibrations-abelian-surfaces-pw}
\source{hitchin-fibrations-abelian-surfaces-pw}
This chapter reproduces the corresponding section of the retrieved
LaTeX source, including its definitions, statements, examples, and
proofs. It has not yet been translated into Lean.
\notready
\end{remark}

% The source section title was: Splittings of perverse filtrations
\subsection{Overview}
In this section, we study splittings of the perverse filtration associated with a proper surjective morphism $\pi: X\to Y$ with $X$ and $Y$ nonsingular. As an application, we strengthen Theorem \ref{taut_abelian} by requiring that the decomposition  {given by said theorem}
\begin{equation}\label{decomp6}
\source{hitchin-fibrations-abelian-surfaces-pw}
H^\ast(\CM_{\beta, A}, \BC) = \bigoplus_{i,d} \widetilde{G}_iH^d(\CM_{\beta, A}, \BC).
\end{equation}
is induced by a ``Lefschetz class" via the mechanisms introduced in \cite{D, dC, dCM4}; see Theorem \ref{strengthened2.1}. As discussed in Remark \ref{remark1}, this is crucial in the study of the specialization morphism (\ref{sp}) in Section 4. 

Throughout, we work with $\BQ$-coefficients except for Section \ref{Section3.5}, where we need to switch to $\BC$-coefficients in order to apply results in Section \ref{Section2.7}. However, we note that all discussions in 
Sections \ref{Lefschetz classes}--\ref{Section3.4} remain valid if we replace $\BQ$ by $\BC$.

\subsection{Lefschetz classes}\label{Lefschetz classes}
\source{hitchin-fibrations-abelian-surfaces-pw}
We consider the perverse filtration associated with $\pi: X \to Y$,
\begin{equation}\label{e50}
\source{hitchin-fibrations-abelian-surfaces-pw}
P_0H^\ast(X, \BQ) \subset P_1H^\ast(X, \BQ) \subset \dots \subset P_{2{R}}H^\ast(X, \BQ) =  H^\ast(X, \BQ),
\end{equation}
where ${R}$ is the defect of semismallness  of $\pi$ (\ref{defect}). The action of a class $\eta \in H^2(X, \BQ)$ on the cohomology $H^*(X, \BQ)$ via cup product satisfies
\begin{equation}\label{e51}
\source{hitchin-fibrations-abelian-surfaces-pw}
\eta: P_kH^i(X, \BQ) \to P_{k+2}H^{i+2}(X, \BQ),
\end{equation}
which further induces an action on 
%\[
%\BH = \bigoplus_{i,j\geq 0} \mathrm{Gr}_i^P H^{i+j}(X, \BQ)= \bigoplus_i \left( P_iH^{i+j}(X, \BQ)/P_{i-1}H^{i+j}(X, \BQ) \right).
%\]
\[
\BH = \bigoplus_{\star,\bullet\geq 0} \mathrm{Gr}_\star^P H^{\bullet}(X, \BQ)= \bigoplus_{\star,\bullet \geq 0} \left( P_\star H^{\bullet}(X, \BQ)/P_{\star-1}H^{\bullet}(X, \BQ) \right).
\]


 We say that  $\eta \in H^2(X, \BQ)$  is a \emph{$\pi$-Lefschetz class}  if its induced action on $\BH$ satisfies the hard Lefschetz-type condition in the sense of \cite[Assumption 2.3.1]{dC}, \emph{i.e.}, the actions
\begin{equation}\label{pi-lefschetz}
\source{hitchin-fibrations-abelian-surfaces-pw}
\eta^k: \mathrm{Gr}^P_{{R} -k}H^*(X, \BQ) \xrightarrow{\simeq}  \mathrm{Gr}^P_{{R}+k}H^*(X, \BQ), \quad \forall k\geq 0,
\end{equation}
are isomorphisms. As a typical example, the relative Hard Lefschetz Theorem \cite{BBD} with respect to $\pi: X \to Y$ implies that a relatively ample class for $\pi$  is a $\pi$-Lefschetz class. The following lemma gives examples of Lefschetz classes other than relatively ample classes.  

\begin{lem}
\source{hitchin-fibrations-abelian-surfaces-pw}
\notready\label{Lemma3.1}
\source{hitchin-fibrations-abelian-surfaces-pw}
Let $A$ be an abelian variety of dimension $m$. Any class $\eta \in H^2(A, \BQ)$ satisfying
\begin{equation}\label{eta1}
\source{hitchin-fibrations-abelian-surfaces-pw}
\eta^m \neq 0 
\end{equation}
is a Lefschetz class with respect to $\pi: A \to \mathrm{pt}$.
\end{lem}

\begin{proof} 
Let $W = H^1(A, \BQ)$. We identify $H^k(A, \BQ)$ with $\wedge^k W$, and therefore $\eta \in \wedge^2W$. The condition (\ref{eta1}) implies that $\eta$ defines a (constant) symplectic form on $W^*$. In particular, we can write $\eta = \sum_{i=1}^m e_i \wedge f_i$ under a basis $e_1,\dots, e_m, f_1, \dots, f_m$ of $W$. Hence the pair $(W\otimes \BC, \eta)$ is the tensor product of $m$ copies of the symplectic plane $(\BC^2, e \wedge f)$, for which the statement is clear.
\end{proof}

Recall the Hitchin fibration $h: \CM_{\mathrm{Dol}} \to \Lambda$. The following proposition is a numerical criterion for Lefschetz classes with respect to $h$.

\begin{prop}
\source{hitchin-fibrations-abelian-surfaces-pw}
\notready\label{prop3.2_Hitchin} Let $F$ be a closed fiber of $h$. A class $\eta \in H^2(\CM_{\mathrm{Dol}}, \BQ)$ is a Lefschetz class with respect to $h$ if and only if {the  following cap product with the fundamental class of the connected fiber $F$ (counting components with multiplicities) is non-trivial}:
\source{hitchin-fibrations-abelian-surfaces-pw}
\begin{equation}\label{eqn50}
\source{hitchin-fibrations-abelian-surfaces-pw}
0 \neq \eta^{\mathrm{dim(F)}} \cap [F] \in H_0(F,\BQ) \simeq \BQ.
%\eta^{\mathrm{dim(F)}}|_{F} \neq 0.
\footnote{
We note that the condition (\ref{eqn50}) does not depend on the choice of the  fiber $F$. }
\end{equation}
\end{prop}



\begin{proof}
 We denote by $\hat{\CM}_{\mathrm{Dol}}$ the corresponding moduli space of stable $\mathrm{PGL}_r$-Higgs bundles with
\[
\hat{h}: \hat{\CM}_{\mathrm{Dol}} \to \hat{\Lambda}
\]
the Hitchin fibration. By \cite{HT1, Markman}, we have
\[
H^2( \hat{\CM}_{\mathrm{Dol}}, \BQ ) = \BQ \cdot \alpha
\]
where $\alpha$ is the first Chern class of an $\hat{h}$-ample line bundle on   $\hat{\CM}_{\mathrm{Dol}}$, and therefore is
a $\hat{h}$-Lefschetz class. By the discussion in \cite[Section 2.4]{dCHM1}, we have an isomorpism
\begin{equation}\label{112233}
\source{hitchin-fibrations-abelian-surfaces-pw}
H^*(\CM_{\mathrm{Dol}}, \BQ) = H^*( \hat{\CM}_{\mathrm{Dol}}, \BQ ) \otimes H^*(\mathrm{Pic}^0(C), \BQ)
\end{equation}
satisfying that
\begin{equation}\label{eqn71}
\source{hitchin-fibrations-abelian-surfaces-pw}
P_k H^*(\CM_{\mathrm{Dol}}, \BQ) = \bigoplus_{i+j=k} P_iH^*( \hat{\CM}_{\mathrm{Dol}}, \BQ ) \otimes H^j(\mathrm{Pic}^0(C), \BQ).
\end{equation}
Under the identification 
\[
H^2(\CM_{\mathrm{Dol}}, \BQ) = H^2( \hat{\CM}_{\mathrm{Dol}}, \BQ ) \oplus H^2(\mathrm{Pic}^0(C), \BQ) = \BQ \cdot \alpha  \oplus H^2(\mathrm{Pic}^0(C)),
\]
induced by (\ref{112233}), we can express any class $\eta \in H^2(\CM_{\mathrm{Dol}}, \BQ)$ as
\begin{equation}\label{123eta}
\source{hitchin-fibrations-abelian-surfaces-pw}
\eta = \mu {\alpha  \otimes 1+ 1\otimes \xi}, \quad \mu \in \BQ,~ \xi \in H^2(\mathrm{Pic}^0(C), \BQ).  
\end{equation}
Lemma \ref{Lemma3.1} {combined with \cite[Appendix]{dC}} implies that the class $\eta$ is Lefschetz if and only if 
\begin{equation}\label{eqn61}
\source{hitchin-fibrations-abelian-surfaces-pw}
\mu \neq 0, \quad \xi^g\neq 0.
\end{equation}
Therefore, it suffices to show that (\ref{eqn61}) is equivalent to the condition (\ref{eqn50}).

We  first consider traceless Higgs bundles
\[
\CM^0_{\mathrm{Dol}} = \{(\CE, \theta) \in \CM_{\mathrm{Dol}}: \mathrm{trace}(\theta)=0\} \subset \CM_{\mathrm{Dol}}.
\]
Recall from \cite[Proposition 2.4.3]{dCHM1} that there is a finite morphism
\[
q: \check{\CM}_{\mathrm{Dol}} \times \mathrm{Pic}^0(C)   \to \CM^0_{\mathrm{Dol}}
\]
with $\check{\CM}_{\mathrm{Dol}}$ the corresponding Dolbeault  moduli  space of stable $\mathrm{SL}_r$-Higgs bundles. The preimage of a closed fiber $F \subset \CM^0_{\mathrm{Dol}}$ of the restricted Hitchin fibration $h|_{\CM^0_{\mathrm{Dol}}}$ is the product
\[
q^{-1}(F) =  \check{F} \times \mathrm{Pic}^0(C),
\]
where $\check{F}$ {is the corresponding  closed fiber} of the $\mathrm{SL}_n$ Hitchin fibration. The pullback of a class (\ref{123eta}) along
\[
\check{F} \times \mathrm{Pic}^0(C) \hookrightarrow \check{\CM}_{\mathrm{Dol}} \times \mathrm{Pic}^0(C)   \xrightarrow{q} \CM^0 \hookrightarrow \CM_{\mathrm{Dol}}
\]
is of the type 
\[
\check{\eta} = \mu \check{\alpha}\otimes 1 + 1\otimes \xi \in H^2(\check{F}, \BQ) \oplus H^2(\mathrm{Pic}^0(C), \BQ) \subset H^2(\check{F} \times \mathrm{Pic}^0(C), \BQ)
\]
with $\check{\alpha}$ an ample class on $\check{F}$. Since $q: \check{F} \times \mathrm{Pic}^0(C) \to F$ is finite
{and surjective}, the condition (\ref{eqn50}) is equivalent to
\[
0 \neq \check{\eta}^{\mathrm{dim}(\check{F} \times \mathrm{Pic}^0(C))} \cap [\check{F}] = \check{\eta}^{\mathrm{dim}(F)}\cap [\check{F}],
\]
which, in turn,  is  equivalent to (\ref{eqn61}). This completes the proof.
\end{proof}

\subsection{The first Deligne splitting}\label{First Deligne Splitting}
\source{hitchin-fibrations-abelian-surfaces-pw}
%Deligne defined in \cite{D} three splittings of the perverse filtration associated with $\pi: X\to Y$.\footnote{Without the assumption that $X$ is nonsingular, the discussion in this section works identically for splitting the perverse filtration on the intersection cohomology $\mathrm{IH}^*(X, \BQ)$.} Each splitting relies on the choice of a $\pi$-Lefschetz class
%\[
%\eta \in H^2(X, \BQ).
%\]
%Deligne's constructions were later extended and generalized in \cite{dC, dCM4}. In this section, we only consider the splitting given by the first construction of \cite{D} which we call \emph{the first Deligne splitting}. 

Given a projective morphism $\pi : X \to Y$ and a $\pi$-ample line bundle  on $X$, Deligne \cite{D} 
constructs three splittings of the direct image $R\pi_* \BQ_X$ (resp. $R\pi_* \mathrm{IC}_X$ if $X$ is singular), which induce splittings of the perverse filtration on the (resp. intersection) cohomology groups $H^*(X,\BQ)$. In this paper, we need the variant \cite{dC} of this construction
where one starts with a $\pi$-Lefschetz class $\eta \in H^2(X,\BQ)$, \emph{i.e.} one that does not necessarily satisfy the relative Hard Lefschetz Theorem in the derived category $D^b_c(Y)$, but nonetheless satisfies the cohomological version (\ref{pi-lefschetz}).
We need only the first of the resulting three splittings, which we name the \emph{first Deligne splitting}
(\cite{D,dCM4,dC}).


%We follow closely the treatment of \cite{dCM4}.


According to \cite{dC}, the first Deligne splitting of the perverse filtration (\ref{e50})   
\[
H^*(X,\BQ) = \bigoplus_{i}G_iH^*(X, \BQ)
\]
 associated with the $\pi$-Lefschetz class $\eta$ (\emph{i.e.} (\ref{pi-lefschetz}) holds)
can be described using \emph{only} the action of $\eta$ on $H^*(X,\BQ)$.  We explain this description explicitly as follows.

For $i  \geq 0$, we let $\mathrm{Gr}^P_i:=P_iH^*(X,\BQ)/P_{i-1}H^*(X,\BQ)$ denote the graded spaces, which are subquotients of cohomology, and let $G_i:= G_i H^*(X,\BQ)$ denote the corresponding image of $\mathrm{Gr}^P_i$ via the first Deligne splitting, i.e. these are the summands in the last paragraph, which are subspaces of cohomology that we want to characterize.


For $0 \leq k \leq {R}$, let 
\[
\mathrm{Gr}^P_k\supseteq Q'_k:= 
\ker \{\eta^{{R}-k+1}: \mathrm{Gr}^P_{k}\to \mathrm{Gr}^P_{2{R} -k +2}\}
\]
be the associated graded-primitive spaces (here $\eta$ acts on the graded spaces of the perverse filtration on cohomology).
Let $Q_k \subseteq G_k$ be the image of $Q'_k$ via the first Deligne splitting. We have
$G_k= \sum_{i \geq 0} \eta^i Q_{k-2i}$ for $0 \leq k \leq {R}$, and $G_{{R}+\kappa}=\eta^k G_{{R}-\kappa}$
(here $\eta$ acts on cohomology).

It follows that, in order to have a complete description of the first Deligne splitting that involves solely
the action of $\eta$ in cohomology, we only need to describe $Q_k \subseteq H^*(X,\BQ)$ in such terms.

We describe $Q_k$ following \cite[Section 2.7]{dCM4}. Note that its context is the one of the fist Deligne splitting arising from working in the derived category, but the proof works verbatim in the present context of cohomology acted upon by  a $\pi$-Lefschetz class. By (\ref{e51}), we have
\[
\eta^j P_k H^*(X, \BQ) \subset P_{k+2j}H^*(X, \BQ).
\]
Let $\Phi_0 (\eta)$ be the composition of the morphisms
\[
\Phi_0 (\eta) : P_kH^*(X, \BQ) \xrightarrow{\eta^{{ R}-k+1}} P_{2{ R}-k+2}H^*(X, \BQ) \to \mathrm{Gr}_{2{ R}-k+2}^PH^*(X, \BQ),
\]
where the first map is cup product and the second map is the projection to the graded piece. We obtain the sub-vector space
\[
\mathrm{Ker}(\Phi_0 (\eta)) \subset P_kH^*(X, \BQ).
\]
The morphisms $\Phi_m (\eta)$ are defined inductively for $m\geq 0$ as follows:
\[
\Phi_m (\eta) : \mathrm{Ker}(\Phi_{m-1} (\eta)) \xrightarrow{\eta^{{ R} -k+m}} P_{2{R} -k+2m}H^*(X, \BQ) \to \mathrm{Gr}_{2 { R} -k+2m}^PH^*(X, \BQ).
\]
Therefore for fixed $k$, we obtain a sequence of sub-vector spaces
\[
\cdots \subset \mathrm{Ker}(\Phi_1(\eta)) \subset \mathrm{Ker}(\Phi_0 (\eta) ) \subset P_kH^*(X, \BQ). 
\]
According to \cite[Proposition 2.7.1]{dCM4}, we have the desired description
\[
Q_k = Q_k (\eta)= \mathrm{Ker}(\Phi_k (\eta)) \subset P_kH^*(X, \BQ), \quad \forall 0 \leq k \leq { R}.
\]



The description of the first Deligne splitting yields immediately the following comparison lemma, which expresses a kind of functoriality for the first Deligne splitting.

{
\begin{lem}
\source{hitchin-fibrations-abelian-surfaces-pw}
\notready\label{comparison1}
\source{hitchin-fibrations-abelian-surfaces-pw}
Let $X_i$ $(i=1,2)$ be nonsingular varieties with proper surjective morphisms $f_i: X_i \to Y_i$, and let $P_\star H^*(X_i, \BQ)$ be the corresponding perverse filtrations.  Let  
\[
\phi: H^*(X_1, \BQ) \to H^*(X_2, \BQ)
\]
be a morphism
of graded $\BQ$-algebras satisfying 
\begin{enumerate}
    \item[(a)] $\phi(P_kH^d(X_1, \BQ)) \subset P_kH^d(X_2, \BQ)$;
    \item[(b)] $\phi(\eta_1) = \eta_2$,  with {$\eta_i \in H^2(X_i, \BQ)$ an $f_i$-Lefschetz class for $i= 1,2$.}
\end{enumerate}
Then we have
\[
\phi(G_kH^*(X_1, \BQ)) \subset G_k H^*(X_2, \BQ), \quad \forall k \geq 0
\]
with $ H^*(X_i,\BQ)= \oplus_{k\geq 0}G_kH^*(X_i, \BQ)$ the first Deligne splitting associated with $\eta_i$, $i=1,2$.
\end{lem}

\begin{proof}
It follows from the description of the first Deligne splittings 
\[H^*(X_i, \BQ) = \bigoplus_{0 \leq k \leq r} 
 \bigoplus_{j \geq 0} 
 \left(\eta^j_i Q_{k-2j}(\eta_i)  \bigoplus \eta^{r-k+j}_i Q_{k-2j}  (\eta_i)
 \right) \quad i=1,2
\] 
summarized above, and the fact that the sub-vector spaces
\[
\mathrm{Ker}(\Phi_k (\eta_i)) = Q_k (\eta_i) \subset P_kH^*(X_i, \BQ)
\]
are preserved under the $P$-filtered morphism $\phi$.
\end{proof}
}

\begin{rmk}
\source{hitchin-fibrations-abelian-surfaces-pw}
\notready
Since cohomologically, the Hitchin fibration 
\[
h: \CM_{\mathrm{Dol}} \to  \Lambda,
\]
behaves like the product of the fibrations $\hat{h}: \hat{\CM}_{\mathrm{Dol}} \to \hat{\Lambda}$ and $\mathrm{Pic}^0(C) \to \mathrm{pt}$ via the ring isomorphism
\[
H^*(\CM_{\mathrm{Dol}}, \BQ) =H^*(\hat{\CM}_{\mathrm{Dol}}, \BQ) \otimes H^*(\mathrm{Pic}^0(C), \BQ)
\]
given in \cite[Section 2.4]{dCHM1}, we see from the proof of Proposition \ref{prop3.2_Hitchin}
 {(cf. (\ref{eqn61}))} together with \cite[Appendix]{dC} that the first Deligne splitting associated with any 
$h$-Lefschetz class, i.e. a class of the form $\mu \alpha\otimes 1+ 1\otimes \xi$ with $\mu\neq 0$ and $\xi^g \neq 0$, has the form
\begin{equation}\label{Hitchin_splitting}
\source{hitchin-fibrations-abelian-surfaces-pw}
G_k H^*(\CM_{\mathrm{Dol}}, \BQ) = \bigoplus_{i+j=k} G_iH^*(\hat{\CM}_{\mathrm{Dol}}, \BQ) \otimes H^j(\mathrm{Pic}^0(C), \BQ), 
\end{equation}
where $G_iH^*(\hat{\CM}_{\mathrm{Dol}}, \BQ)$ is the first Deligne splitting associated with the $\hat{h}$-Lefschetz class 
\[
\alpha \in H^2(\hat{\CM}_{\mathrm{Dol}}, \BQ).
\]
\begin{rmk}
\source{hitchin-fibrations-abelian-surfaces-pw}
\notready\label{dolglcansplit}
\source{hitchin-fibrations-abelian-surfaces-pw}
In particular, the splitting (\ref{Hitchin_splitting}) does not depend on the choice of an  $h$-Lefschetz class.
\end{rmk}
For  every genus $g \geq 2$, we expect, and actually prove in the $g=2$ case, that (\ref{Hitchin_splitting}) serves as the splitting in Conjecture \ref{P=W2}.
\end{rmk}


\subsection{Semismall maps and Hilbert schemes}\label{Section3.4}
\source{hitchin-fibrations-abelian-surfaces-pw}
In this section, we study the situation where our morphism $\pi: X \to Y$ can be factored as
\[
X \xrightarrow{f} Z \xrightarrow{g} Y
\]
with $f: X \to Z$ semismall and surjective; in particular, $f$ is generically finite and $\dim (X)=\dim (Z)$. We further assume that there are closed irreducible sub-varieties $Z_i \subset Z$ such that the decomposition theorem for $f$ takes  the form of a canonical  finite direct sum decomposition
\begin{equation}\label{decomp9}
\source{hitchin-fibrations-abelian-surfaces-pw}
Rf_* \BQ_X = \bigoplus_i \mathrm{IC}_{Z_i}[-\mathrm{dim}(X)],
\end{equation}
where each $\mathrm{IC}_{Z_i}$ is the (perverse) intersection cohomology complex of $Z_i$. Note that we have the following identity concerning intersection cohomology groups
$\mathrm{IH}^d(Z_i, \BQ)=H^{d-\dim Z_i}(Z_i, \mathrm{IC}_{Z_i})$.
One of the $Z_i$ is the total variety $Z$. The restriction of $g$ to each $Z_i$ yields the morphism
\[
g_i: Z_i \to Y_i \subset Y.
\]
We deduce from (\ref{decomp9}) a canonical decomposition of the cohomology of $X$
\begin{equation}\label{decomp8}
\source{hitchin-fibrations-abelian-surfaces-pw}
H^d(X, \BQ) = \bigoplus_i \mathrm{IH}^{d-c_i}(Z_i, \BQ), {\quad c_i = \dim X - \dim Z_i={\rm codim}_X Z_i.}
\end{equation}

Let ${R}$ be the defect of semismallness of $\pi$ and, for each $i$, let ${R_i}$ be the defect of semismallness of $g_i:Z_i \to Y_i$.

\begin{prop}
\source{hitchin-fibrations-abelian-surfaces-pw}
\notready\label{proposition3.5}
\source{hitchin-fibrations-abelian-surfaces-pw}
Let  $\alpha \in H^2(Z, \BQ)$ satisfy that, for every $i$,   the restriction
\[
\alpha_i = \alpha|_{Z_i} \in H^2(Z_i, \BQ)
\]
is a $g_i$-Lefschetz class  with associated  first Deligne splitting 
\[
\mathrm{IH}^*(Z_i, \BQ) = \bigoplus_k G_k \mathrm{IH}^*(Z_i, \BQ).
\]
Then $\eta = f^*\alpha \in H^2(X, \BQ)$ is a $\pi$-Lefschetz class, whose associated first Deligne splitting is,
under the identification (\ref{decomp8}), given by
\begin{equation}\label{decomp7}
\source{hitchin-fibrations-abelian-surfaces-pw}
{G_kH^d(X, \BQ) = \bigoplus_i G_{k-R+{R_i}}\mathrm{IH}^{d-c_i}(Z_i, \BQ).}
\end{equation}
\end{prop}

\begin{proof}{Recall that we have defined the perverse filtration $P$ on $H^*(X,\BQ)$ concentrated in the interval $[0,{2R}]$,
and similarly,  for every $i$, the perverse filtration  $P$ on $\rm{IH}^*(Z_i,\BQ)$ is concentrated in the interval
$[0,{2R_i}]$. It follows that, according to  (\ref{decomp9}) and (\ref{decomp8}),  the direct summand $\rm{IH}^*(Z_i,\BQ)$ contributes to  
$H^*(X,\BQ)$ in perversities in the interval ${[R-R_i,R+R_i]}$.

We apply the decomposition theorem to the composition
\[
X \to Z \to Y
\]
and obtain that the perverse filtration $P_\star H^*(X, \BQ)$ can be expressed in terms of the perverse filtrations $P_\star \mathrm{IH}^*(Z_i, \BQ)$ under (\ref{decomp8}), \emph{i.e.},
\[
P_kH^d(X, \BQ) = \bigoplus_i P_{k-{R} +{R_i}}\mathrm{IH}^{d-c_i}(Z_i, \BQ).
\]


By \cite[Remark 4.4.3]{dCM0}, the action of $\eta=f^*\alpha$ on the l.h.s. of (\ref{decomp8}), is 
the direct sum of the actions of the classes $\alpha_i=\alpha_{|Z_i}$ on the summands of the r.h.s..


Since every $\alpha_i$ is a $g_i$-Lefschetz class, 
we deduce that $\eta$ is a $\pi$-Lefschetz class with associated first Deligne splitting decomposition (\ref{decomp7}).}
\end{proof}


{Next, we show that, given a fibered abelian surface 
\[p:A=E\times E' \to E'\]
  as in Section \ref{Section1.5}, 
the splitting (\ref{canonical_split}) of the perverse filtration on $H^*(A^{[n]}, \BQ)$ associated with the natural morphism  $p_n: A^{[n]} \to E'^{(n)}$, 
 is the first Deligne splitting induced by a  $p_n$-Lefschetz class. 

The morphism $p_n: A^{[n]} \to E'^{(n)}$ admits the natural factorization 
\begin{equation}\label{factorization_Hilb}
\source{hitchin-fibrations-abelian-surfaces-pw}
A^{[n]} \xrightarrow{f} A^{(n)} \xrightarrow{g} E'^{(n)}
\end{equation}
where the Hilbert--Chow morphism $f$ is semismall \cite{Go2}. There are canonical morphisms
\[
\kappa_\nu: A^{(\nu)} \to A^{(n)}
\]
together with a canonical stratification of $A^{(n)}$ indexed by the partitions $\nu$ of $n$,
\[
A^{(n)} = \bigcup_\nu A_\nu, \quad \overline{A_\nu} = \mathrm{Im}(\kappa_\nu) \subset A^{(n)}.
\]
Note that  the resulting morphism $\kappa_\nu:A^{(\nu)} \to \overline{A_\nu}$ is the normalization of the target,
so that 
\begin{equation}\label{hisih}
\source{hitchin-fibrations-abelian-surfaces-pw}
H^*(A^{(\nu)},\BQ)= {\rm IH}^*(\overline{A_\nu},\BQ).
\end{equation}

By \cite[Theorem 3]{Go2}, the decomposition theorem associated with $f$ takes the form
analogous to (\ref{decomp9}),
\[
Rf_* \BQ_{A^{[n]}} = \bigoplus_{\nu} \mathrm{IC}_{\overline{A_\nu}} [-2n],
\]
and the restriction of $g: A^{(n)} \to E'^{(n)}$ to each $\overline{A_\nu}$ is described by the commutative diagrams
\begin{equation}\label{diagram9}
\source{hitchin-fibrations-abelian-surfaces-pw}
\begin{tikzcd}
A^{(\nu)} \arrow{r}{\kappa_\nu}  \arrow[swap]{d} &  \overline{A_\nu}\arrow[r, hookrightarrow]  \arrow[swap]{d}{p_\nu} &  A^{(n)}  \arrow{d} \\
{E'}^{(\nu)} \arrow{r} & \overline{{E'_\nu}} \arrow[r, hookrightarrow] & E'^{(n)}.
\end{tikzcd}
\end{equation}

We consider the $p$-ample class
\[
\alpha = \mathrm{pt} \otimes 1_{E'} \in H^2(E, \BQ)\otimes H^0(E', \BQ) \subset H^2(A, \BQ)
\]
of $p:A \to E'$ which induces a Lefschetz class 
\begin{equation}\label{alpha(n)}
\source{hitchin-fibrations-abelian-surfaces-pw}
\alpha^{(n)} \in H^2(A^{(n)}, \BQ)
\end{equation}
with respect to $A^{(n)} \to E'^{(n)}$. It further induces a Lefschetz class
\[
\alpha^{(\nu)} \in H^2(A^{(\nu)}, \BQ)
\]
with respect to $A^{(\nu)} \to E'^{(\nu)}$ for every partition $\nu$ of $n$. By the diagram (\ref{diagram9}), the pullback of $\alpha^{(n)}$ to every $A^{(\nu)}$ via $\kappa_\nu$ coincides with $\alpha^{(\nu)}$.

By keeping in mind that  ${\rm codim}_{A^{(n)}} \overline{A_\nu}=2n -2l(\nu)$, that the defects of semismallness of the morphisms  $p_n$ and $p_\nu$ are $n$ and $l(\nu)$, respectively, and the identity (\ref{hisih}),
we obtain the following corollary by applying Proposition \ref{proposition3.5} to the factorization (\ref{factorization_Hilb}).
}

\begin{cor}
\source{hitchin-fibrations-abelian-surfaces-pw}
\notready\label{cor3.6}
\source{hitchin-fibrations-abelian-surfaces-pw}
The decomposition (\ref{canonical_split}) is the first Deligne splitting induced by the 
$(p_n: A^{[n]} \to E'^{(n)})$-Lefschetz class  ($f$ as in (\ref{factorization_Hilb}), $\alpha^{(n)}$
as in (\ref{alpha(n)}))
\[
\eta_{A^{[n]}} = f^*\alpha^{(n)} \in H^2(A^{[n]}, \BQ).
\]
\end{cor}

\subsection{A strengthened version of Theorem \ref{taut_abelian}}\label{Section3.5}
\source{hitchin-fibrations-abelian-surfaces-pw}
We study the decomposition (\ref{decomp6}) constructed for Theorem \ref{taut_abelian}. In the special case (\ref{special}), we have
\[
\CM_{\beta, A} \cong A^{[n]} \times \hat{A}, \qquad A=E \times E',
\]
with the morphism $\pi_\beta: \CM_{A,\beta} \to B$ given by the morphism $p_n\times q$ in (\ref{HC2}). Corollary \ref{cor3.6} 
implies that the decomposition (\ref{splitting2}) is the first Deligne splitting associated with the 
$(p_n\times q)$-Lefschetz class 
\[
\eta_{A,\beta} = \eta_{A^{[n]}} \boxtimes 1_{\hat{A}} + 1_{A^{[n]}}\boxtimes \left(1_E \boxtimes \mathrm{pt} \right) \in H^2(A^{[n]} \times \hat{A}, \BQ) =  H^2(\CM_{\beta, A}, \BQ).
\]

By Theorem \ref{thm2.6}, a pair $(A', \beta')$ with $\beta'^2=2n$ is perversely isomorphic to the special pair $(A, \beta)$ given by (\ref{special}). So there is a graded isomorphism
\[
\phi: H^*(\CM_{\beta, A}, \BC) \xrightarrow{\simeq} H^*(\CM_{\beta', A'}, \BC) 
\]
of $\BC$-algebras preserving the corresponding perverse filtrations, and we obtain a decomposition (\ref{decomp6}) as
\begin{equation}\label{996}
\source{hitchin-fibrations-abelian-surfaces-pw}
\widetilde{G}_kH^d(\CM_{\beta',A'}, \BC) = \phi(\widetilde{G}_kH^d(\CM_{\beta',A'}, \BC));
\end{equation}
see the proof of Proposition \ref{prop2.5}. Hence
\[
\eta_{\beta',A'} = \phi(\eta_{\beta,A}) \in H^2(\CM_{\beta',A'}, \BC)
\]
is a $\pi_{\beta'}$-Lefschetz class, and Lemma \ref{comparison1} implies that the decomposition (\ref{996}) is the first Deligne splitting associated with $\eta_{\beta',A'}$. 

This gives the following strengthened version of Theorem \ref{taut_abelian}, which, we stress, is about any pair $(A,  \beta)$ with $v_\beta = (0, \beta, \chi)$ primitive, not just the special cases (\ref{special}).

\begin{thm}
\source{hitchin-fibrations-abelian-surfaces-pw}
\notready\label{strengthened2.1}
\source{hitchin-fibrations-abelian-surfaces-pw}
Theorem \ref{taut_abelian} holds for a splitting $\widetilde{G}_\ast H^\ast(\CM_{\beta,A}, \BC)$ given by the first Deligne splitting associated with a {suitable $\pi_\beta$-Lefschetz class $\eta_\beta$}, where $\pi_\beta: \CM_{\beta,A} \to B$  is the morphism (\ref{Hilb-Ch}).
\end{thm}
